% Test of margin floats (didactic) in the notes: memoir + beamerarticle,
% the source in test-marginfloat-contents.tex shared with
% test-marginfloat-slides.tex.
%
% Build: see the Makefile (twice, for the references and
% marginfix's anchors).
%
% Expected in the resulting PDF, all on page 1:
%   1  The ordinary figure PENDINGCAP (Figure 1) at the bottom of the page.
%      It is waiting when the margin floats come, so memoir's own
%      environments would start a new page there (\FloatBlock); ours
%      don't, so the whole text is on page 1 and the PDF has one page.
%   2  The margin figure (example-image-a, 3cm high, "A picture beside its
%      text.", Figure 2) in the right margin, beside the paragraph that
%      starts "Figure 2 is in the margin", not in the text block.  It is
%      pushed down by the note NOTEABOVE just above it and doesn't touch
%      it: a gap of \marginparpush between the note's last line and the
%      image.
%   3  The margin table (MTABCAP, Table 1) in the right margin, as close
%      to its line (SPACEA SPACEB) as the figure above lets it, set ragged
%      right; TABCELL is not indented, but flush with the caption.
%   4  One space between SPACEA and SPACEB, as between other words.
%   5  REFS: "Figure 1, Figure 2 and Table 1", with no undefined
%      references in the log.
\documentclass[a4paper]{memoir}
\usepackage[british]{babel}
\usepackage{beamerarticle}
\usepackage{graphicx}
\usepackage{lipsum}
\usepackage{hyperref}
\usepackage[capitalise,noabbrev]{cleveref}
\usepackage{didactic}
\counterwithout{figure}{chapter}
\counterwithout{table}{chapter}
\begin{document}
% The source shared by test-marginfloat.tex (the notes) and
% test-marginfloat-slides.tex (the slides): what a deck writes.  In beamer,
% the text between the frames is skipped.
\mode*
\begin{frame}
  \frametitle{A float that waits}
  PENDINGANCHOR This frame has an ordinary figure, placed at the bottom
  of the page, so it is still waiting when the margin floats come.
  \begin{figure}[b]
    \centering
    \rule{0.5\linewidth}{1.5cm}
    \caption{PENDINGCAP, an ordinary figure.}\label{fig:pending}
  \end{figure}
\end{frame}

\lipsum[1][1-6] The note is on this line,
NOTEANCHOR\ltnote{NOTEABOVE \lipsum[2][1-3]}
just above the paragraph with the figure.

\mode<all>
\begin{frame}
  \begin{marginfigure}
    \includegraphics[width=\linewidth, height=3cm,
      keepaspectratio]{example-image-a}
    \caption{A picture beside its text.}\label{fig:beside}
  \end{marginfigure}
  \Cref{fig:beside} is in the margin of the notes.
\end{frame}

\mode*

\begin{frame}
  \frametitle{A table}
  The table stands in the paragraph: SPACEA \begin{margintable}
    \caption{MTABCAP, a table in the margin.}\label{tab:margin}
    \begin{tabular}{lr}
      TABCELL & 1 \\
      two & 2 \\
    \end{tabular}
  \end{margintable} SPACEB, with one space between the words.
  \lipsum[3][1-4]
\end{frame}

\begin{frame}
  \frametitle{References}
  REFS \cref{fig:pending}, \cref{fig:beside} and \cref{tab:margin}.
\end{frame}

\end{document}
